\begin{frame}[t]{1.3 O que é um paradigma?}
Um \textbf{paradigma} é uma forma de organizar o raciocínio e a construção de programas.
\begin{itemize}
\item Imperativo: indicar comandos que alteram o estado.
\item Procedural: organizar comandos em procedimentos ou funções.
\item Orientado a objetos: agrupar estado e operações em objetos.
\item Funcional: compor funções e transformações de valores.
\end{itemize}
\begin{block}{Linguagens podem combinar paradigmas}
C é usado principalmente de forma imperativa e procedural; C++ permite, entre outras abordagens, programação procedural e orientada a objetos.
\end{block}
\end{frame}
